\chapter{Heat Equation}

\section{Fundamental solution and representation formulas}

We consider the heat equation
$$
u_t-\Delta_x u=0.
$$

\begin{definition}[Fundamental solution of the heat equation]
The fundamental solution of the heat equation on $\mathbb{R}^n$ is the function
$
\Phi:\mathbb{R}^n\times \mathbb{R}\setminus\{0\}\to \mathbb{R}
$
defined by
$$
\Phi(x,t)=
\begin{cases}
\dfrac{1}{(4\pi t)^{n/2}}\exp\!\left(-\dfrac{|x|^2}{4t}\right), & x\in \mathbb{R}^n,\ t>0,\\[1em]
0, & x\in \mathbb{R}^n,\ t<0.
\end{cases}
$$
\end{definition}

\begin{lemma}[Integral of the fundamental solution]
For each $t>0$,
$$
\int_{\mathbb{R}^n}\Phi(x,t)\,dx=1.
$$
\end{lemma}

\begin{theorem}[Solution of the initial-value problem]
Consider the initial value problem
$$
\begin{cases}
u_t-\Delta u=0 & \text{in } \mathbb{R}^n\times (0,\infty),\\
u=g & \text{on } \mathbb{R}^n\times \{t=0\}.
\end{cases}
$$
Let $u(x,t)$ be defined by
$$
u(x,t)=\int_{\mathbb{R}^n}\Phi(x-y,t)g(y)\,dy.
$$
Then:
\begin{enumerate}
    \item $u\in C^\infty(\mathbb{R}^n\times \mathbb{R}_+)$,
    \item $u$ solves the heat equation,
    \item $u$ satisfies the initial value
    $
    \lim_{(x,t)\to (x^0,0)}u(x,t)=g(x^0).
    $
    for $t >0$.
\end{enumerate}
\end{theorem}


\begin{theorem}[Solution of the inhomogeneous problem]
Consider the inhomogeneous problem
$$
\begin{cases}
u_t-\Delta u=f, & (t,x)\in \mathbb{R}_+\times \mathbb{R}^n,\\
u=0, & (t,x)\in \{0\}\times \mathbb{R}^n.
\end{cases}
$$
Then
$$
u(x,t)=\int_0^t\int_{\mathbb{R}^n}\Phi(x-y,t-s)f(y,s)\,dy\,ds.
$$
Moreover:
\begin{enumerate}
    \item $u\in C^\infty$,
    \item $u$ solves the PDE,
    \item
    $
    \lim_{(x,t)\to (x^0,0)}u(x,t)=0.
    $
\end{enumerate}
\end{theorem}

\begin{theorem}[Solution of the inhomogeneous problem with initial data]
Combining the above two theorems, we discover that
$$
u(x,t)
=
\int_{\mathbb{R}^n}\Phi(x-y,t)g(y)\,dy
+
\int_0^t\int_{\mathbb{R}^n}\Phi(x-y,t-s)f(y,s)\,dy\,ds
$$
is a solution of
$$
\begin{cases}
u_t-\Delta u=f & \text{in } \mathbb{R}^n\times (0,\infty),\\
u=g & \text{on } \mathbb{R}^n\times \{t=0\},
\end{cases}
$$
under the hypotheses on $g,f$ as above (i.e.
$
g\in C(\mathbb{R}^n)\cap L^\infty(\mathbb{R}^n), 
$
and 
$
f\in C_c^2(\mathbb{R}^n\times [0,\infty))
$
with compact support.)
\end{theorem}

\section{Mean value formula}

\subsection*{Motivation}
Harmonic functions are characterized by averages over balls, where balls arise as level sets of the fundamental solution of $\Delta u=0$.
For the heat equation, we look at level sets of $\Phi(t-s,x-y)$:
$$
E(x,t;r):=
\left\{
(y,s)\in \mathbb{R}^{n+1}: s\le t,\ \Phi(x-y,t-s)\ge \frac{1}{r^n}
\right\}.
$$
This is called a \emph{heat ball}.

\begin{theorem}[Mean value property for the heat equation]
Let
$$
u\in C_1^2(U_T)
:=
\left\{
u\in C(U_T): D_xu,\ D_x^2u,\ u_t\in C
\right\}.
$$
Then
$$
u(x,t)
=
\frac{1}{4r^n}
\iint_{E(x,t;r)}
u(y,s)\frac{|x-y|^2}{(t-s)^2}\,dy\,ds.
$$
\end{theorem}

\begin{remark}
Here $C_1^2$ means:
\begin{itemize}
    \item ``$2$'': up to second order derivatives in space,
    \item ``$1$'': up to first order derivative in time.
\end{itemize}
\end{remark}

\section{Maximum principle, positivity, uniqueness}

\begin{theorem}[Maximum principle]
Here we have
\begin{enumerate}
    \item
    Weak maximum principle
    $$
    \max_{\overline{U_T}}u=\max_{\Gamma_T}u.
    $$

    \item If $U$ is connected and there exists $(x_0,t_0)\in U_T$ such that
    $$
    u(x_0,t_0)=\max_{\overline{U_T}}u,
    $$
    then $u$ is constant in $\overline{U_{t_0}}$.
    
    This is the strong maximum principle.
\end{enumerate}
\end{theorem}

\begin{theorem}[Positivity]
Let $U$ be connected, and let
$
u\in C_1^2(U_T)\cap C(\overline{U_T})
$
satisfy
$$
\begin{cases}
u_t-\Delta u=0 & \text{in } U_T,\\
u=0 & \text{on } \partial U\times [0,T],\\
u=g & \text{on } U\times \{t=0\}.
\end{cases}
$$
If $g\ge 0$ and there exists $x_0\in U$ such that $g(x_0,0)>0$, then
$
u>0
$
everywhere.
\end{theorem}

\begin{theorem}[Uniqueness on bounded domains]
Let
$
g\in C(\Gamma_T)
$
and
$
f\in C(U_T).
$
Then there exists at most one solution
$
u\in C_1^2(U_T)\cap C(\overline{U_T})
$
of the boundary value problem
$$
\begin{cases}
u_t-\Delta u=f & \text{in } U_T,\\
u=g & \text{on } \Gamma_T.
\end{cases}
$$
\end{theorem}

\subsection*{Now we study $U=\mathbb{R}^n$}

\begin{theorem}[Maximum principle for the Cauchy problem]
Suppose
$
u\in C_1^2(\mathbb{R}^n\times (0,T))\cap C(\mathbb{R}^n\times [0,T])
$
solves
$$
\begin{cases}
u_t-\Delta u=0 & \text{in } \mathbb{R}^n\times (0,T],\\
u=g & \text{on } \mathbb{R}^n\times \{t=0\},
\end{cases}
$$
and satisfies the growth estimate
$$
u(x,t)\le A e^{a|x|^2}
$$
for constants $A,a>0$.
Then
$$
\sup_{\mathbb{R}^n\times [0,T]}u(x,t)=\sup_{\mathbb{R}^n}g(x).
$$
\end{theorem}


\begin{theorem}[Uniqueness for the Cauchy problem]
Let
$
g\in C(\mathbb{R}^n)
$
and
$
f\in C(\mathbb{R}^n\times [0,T]).
$
Then there is at most one solution
$
u\in C_1^2(\mathbb{R}^n\times (0,T))\cap C(\mathbb{R}^n\times [0,T])
$
of the initial value problem
$$
\begin{cases}
u_t-\Delta u=f & \text{in } \mathbb{R}^n\times (0,T],\\
u=g & \text{on } \mathbb{R}^n\times \{t=0\},
\end{cases}
$$
satisfying the growth estimate
$$
|u(x,t)|\le A e^{a|x|^2}
$$
for constants $A,a>0$.
\end{theorem}

\begin{proof}
Apply the theorem on the maximum principle for the Cauchy problem.
\end{proof}

\begin{remark}
Without growth control, this is not true.
\end{remark}

\section{Regularity}

\begin{definition}[Parabolic Cylinder]
    The cylinder is defined as $$C(x,t;r)=\{(y,s)\mid |x-y|\le r,\ t-r^2\le s\le t\}$$
    to denote the closed circular cylinder of radius $r$, height $r^2$, and top center point $(x,t)$.
\end{definition}

\begin{theorem}[Smoothness]
Suppose
$
u\in C_1^2(U_T)
$
solves the heat equation. Then
$
u\in C^\infty(U_T).
$
\end{theorem}

\begin{remark}
This is true even if $u$ is not smooth on $\Gamma_T$.
\end{remark}

\begin{proof}
Leverage smoothness of the convolution formula, together with results on cutoffs and uniqueness.
\end{proof}

\begin{theorem}[Estimate on derivatives]
For all $k,\ell\in \mathbb{N}$, there exists $C_{n,k,\ell}$ such that
$$
\max_{Q(x,t;r/2)}\left|D_x^kD_t^\ell u\right|
\le
\frac{C_{n,k,\ell}}{r^{k+2\ell+n+2}}
\|u\|_{L^1(Q(x,t;r))}
$$
for suitable parabolic cylinders $Q(x,t;r)$.
\end{theorem}

\begin{remark}
The norm on $D_x^kD_t^\ell u$ is the Frobenius norm.
\end{remark}

\section{Energy methods and uniqueness}

\begin{theorem}[Uniqueness]
There is at most one solution of the boundary value problem
$$
\begin{cases}
u_t-\Delta u=f & \text{in } U_T,\\
u=g & \text{on } \Gamma_T,
\end{cases}
$$
with
$
u\in C_1^2(\overline{U_T}).
$
\end{theorem}


\begin{theorem}[Backward uniqueness]
Suppose $u,\hat{u}$ solve
$$
\begin{cases}
u_t-\Delta u=0 & \text{in } U_T,\\
u=g & \text{on } \partial U\times [0,T],
\end{cases}
$$
with
$
u,\hat{u}\in C_1^2(\overline{U_T}).
$
Assume at time $T$ we have $u(x,T)=\tilde{u}(x,T)$ for all $x\in U$.
Then
$
u\equiv \hat{u}
$
within $U_T$.

\end{theorem}
